\chapter{Transport Adapters}
\label{ch:transports}

\begin{note}
Stub. Three transports delegate to \code{GraphHunterApi}: the Tauri
desktop app, the axum-based HTTP server, and the TypeScript MCP
server. All three are meant to be thin --- no logic beyond
constructing a request DTO and rendering the response.
\end{note}

\section{Tauri}

Commands live in \filepath{app/src-tauri/src/commands/*.rs}.
Every command body is $\leq$ 10 LOC:

\begin{lstlisting}[style=rust]
#[tauri::command]
pub fn cmd_run_hunt(
    api: State<Arc<GraphHunterApi>>,
    hypothesis: Hypothesis,
    time_window: Option<(i64, i64)>,
    max_results: Option<u32>,
) -> Result<RunHuntResponse, CommandError> {
    api.run_hunt(RunHuntRequest {
        session: None,
        hypothesis,
        time_window,
        max_results,
    })
    .map_err(CommandError::from)
}
\end{lstlisting}

Commands are registered in \filepath{app/src-tauri/src/lib.rs} with
the \code{tauri::Builder::invoke\_handler} macro.

\section{HTTP}

Handlers live in
\filepath{app/src-tauri/src/http\_api.rs} and share the same
\code{Arc<GraphHunterApi>} via axum's \code{State}. The HTTP port
is chosen at startup (\code{util::resolve\_api\_port}) and written
to the Tauri app's config directory so the MCP server and CLI can
find it.

Every handler is also $\leq$ 10 LOC:

\begin{lstlisting}[style=rust]
async fn handler_run_hunt(
    State(api): State<Arc<GraphHunterApi>>,
    Json(body): Json<RunHuntRequest>,
) -> Response {
    match api.run_hunt(body) {
        Ok(v)  => ok_json(v),
        Err(e) => error_json(e),
    }
}
\end{lstlisting}

\section{MCP}

\filepath{graph-hunter-mcp/src/index.ts} exposes $\sim$30 tools
over the MCP protocol. Each tool handler HTTPs-calls the local
engine via \code{http://127.0.0.1:<port>}; it does not link the
Rust engine at all. This is the only transport that doesn't hold
an \code{Arc<GraphHunterApi>} directly --- because it runs in a
different process (Node.js) --- yet it sees the exact same
semantics because the HTTP layer is a thin wrapper around the
façade.

\section{EventEmitter}

Progress events (session load, ingest, Sentinel streaming) flow
through an \code{Arc<dyn EventEmitter>} stored on \code{ApiState}.
\code{TauriEventEmitter} bridges to \code{AppHandle::emit};
\code{NoopEmitter} is used by HTTP, MCP, CLI, and tests.

\begin{inthecode}
\filepath{app/src-tauri/src/commands/} --- Tauri command modules.\\
\filepath{app/src-tauri/src/http\_api.rs} --- HTTP handlers.\\
\filepath{graph-hunter-mcp/src/index.ts} --- MCP tools.\\
\filepath{graph\_hunter\_api/src/events.rs} (or similar) ---
\code{EventEmitter} trait.
\end{inthecode}
